\section{Experimental design}\label{sec:design}

\paragraph{Instances and budget.} We use Xi'an and Amsterdam with a 20-station budget and the chained reward (Section~\ref{sec:benchmark}). The benchmark's own scripts use longer lines (47 stations for Xi'an and 21 for Amsterdam). We use 20 stations on both grids, and we did not run the longer setting, so our numbers are not comparable with figures published for those settings. Every method evaluates about $10^4$ lines: 10\,000 episodes for the RL policy and tabular Q-learning, 10\,100 lines for NSGA-II and the full-budget scalarised GA, 9\,900 for the matched scalarised GA, and the default 50 generations of population 100 for the GA of \citet{michailidis2026}, which completes about 5\,100 lines because infeasible offspring are discarded. Evaluation seeds are 0--9. Seeds 100--102 are used only for the objective anchors and for hyperparameter selection.

\paragraph{Hyperparameter selection.} For the RL policy we compared six configurations on the tuning seeds at the full budget and chose the one with the highest mean expected utility of the greedy 15-point front. NSGA-II received the same treatment over six population and generation splits at a constant total budget. The scalarised GA uses a fixed split and is not tuned. The analysis plan, including the primary endpoints and the rules for what may be claimed, was written before any evaluation run and is included in the repository.

\paragraph{Deployable set and archive.} Every method yields two sets of lines. The \emph{deployable set} is what a planner would receive for the 15 weight vectors: the greedy line of the trained policy, the best member of the final population for each weight vector, or the single line of a method that has no weight input. The \emph{archive} is the non-dominated subset of all lines the method evaluated during the run. Archive indicators measure search quality, deployable indicators measure what the method hands over.

\paragraph{Indicators.} On the scaled objectives we report the hypervolume with reference point at the origin, which is the value of an empty line, and the expected utility $\mathrm{EUM}=\frac{1}{15}\sum_{\mathbf w}\max_{\tilde{\mathbf f}}\{-g(\tilde{\mathbf f}\mid\mathbf w)\}$ over the 15 weight vectors \citep{hayes2022}. Inverted generational distance plus \citep{ishibuchi2015} is computed against the non-dominated union of all runs. We also report the raw objectives, the worst-served share $s_{\min}$ and the generalised Gini welfare at the three vertex weights and the centroid.

\paragraph{Statistics.} Methods are compared with two-sided Mann--Whitney tests, the Vargha--Delaney $A_{12}$ effect size and 95\% bootstrap intervals of the mean difference (10\,000 resamples). Paired comparisons across conditions of the same seed use the exact Wilcoxon signed-rank test and Hodges--Lehmann estimates. Holm correction is applied within each table. With ten seeds, an exact two-sided Wilcoxon test cannot give $p<0.002$, so after Holm correction over several contrasts only effects present in nearly every seed can be called significant; we state this where it matters.

\paragraph{Ablations.} We remove one component at a time: the multi-objective weights (demand-only reward), message passing, the lookahead features, the learned start policy, the scalarised prior of Section~\ref{sec:method}, the generalised Gini reward (replaced by one minus Gini) and the objective scaling.

\paragraph{Leave-one-line-out test.} We delete Line~2 from the Xi'an instance, keep the exclusions that apply to Line~1, and ask each method for a 27-station extension, the length of the real Line~2. We measure precision, recall and F1 of the generated cells against the real Line~2 within one cell, the mean symmetric distance, the angle between principal axes, and the demand ratio to the real line. Each metric is converted to a percentile among 10\,000 random feasible 27-cell lines. Two limitations apply and are repeated in Section~\ref{sec:discussion}: the OD matrix was recorded when both lines were operating, and Line~2 opened before Line~1.

\paragraph{Cross-city warm start.} Policies trained on one city initialise training on the other in six conditions: scratch, all transferable parameters, encoder only, heads only, all parameters with a frozen encoder, and a random frozen encoder. The last condition separates transferred knowledge from having fewer trainable parameters. Runs are paired by seed and use the full and a reduced (20\%) episode budget.
